% ============================================================================
\section{Does Supervision Density Change Gradient Sparsity?}
\label{sec:density}
% ============================================================================

Our third study asks whether vocabulary-level supervision changes how
strongly gradients concentrate on a small fraction of model parameters.
Sampled-token PG uses feedback for the generated token, full-vocabulary
reverse KL includes every vocabulary alternative, and teacher top-$k$
provides a practical intermediate choice
(Section~\ref{sec:setting}). This concerns vocabulary coverage at each
prefix, rather than the number of supervised response positions.
Sampled-token feedback does not imply a sparse parameter gradient:
the sampled token's log probability depends on all logits through the
softmax normalizer. Concentrated OPD parameter changes have also been
observed with dense feedback \citep{yu2026denseupdates}. We therefore
test for a difference without assuming that PG produces sparser gradients
or better capability integration.

\subsection{Matched single-teacher and joint experiments}

We select one representative single-teacher setting before inspecting
results and compare PG with corrected teacher top-$64$ in that setting
and in joint MOPD (Table~\ref{tab:density_controls}). These four online
configurations use the losses in Equations~\ref{eq:mopd_pg}
and~\ref{eq:mopd_topk}; the joint top-$64$ run also anchors
Section~\ref{sec:normalization}. All use domain-response averaging and
match initialization, teachers, optimizer, loss masks, response cap,
prompt schedule, and update budget within each setting. The top-$64$
loss uses probabilities normalized over the full vocabulary, without
renormalizing within the selected tokens; we record any clipping.
The shared model and training setup appears in
Section~\ref{sec:normalization} and Appendix~\ref{app:setup}.

\begin{table}[!htbp]
\centering
\small
\begin{tabular}{lll}
\toprule
Setting & Online losses & Common-prefix reference \\
\midrule
Single teacher & PG, teacher top-$64$ & Full-vocabulary reverse KL \\
Joint MOPD & PG, teacher top-$64$ & Full-vocabulary reverse KL \\
\bottomrule
\end{tabular}
\caption{Supervision-density experiments. Four online configurations
share full-vocabulary local measurements at matched checkpoints.}
\label{tab:density_controls}
\end{table}

At the initial, early, intermediate, and final checkpoints, we evaluate
all three losses on a small shared bank of student-generated prefixes.
Each comparison holds the student, prefix weights, and optimizer state
fixed; local PG uses fresh token samples at those prefixes.
Full-vocabulary reverse-KL gradients and proposed optimizer steps
are a primary local reference in both settings; a full-vocabulary online
run is conditional on measured resource feasibility. Online PG and
top-$64$ branches generate their own on-policy responses, so their
trajectories compare practical loss choices rather than isolating
vocabulary coverage as the only causal difference.

\subsection{Gradient concentration and learning outcomes}

Our primary measurements use raw gradients before clipping. Following
Section~\ref{sec:subnetworks}, we report the fraction of coordinates below
each magnitude threshold, gradient norms, and the fraction of squared
gradient magnitude carried by the largest coordinates. Repeated batches
on the common prefix bank quantify sampling variability. We separately
measure proposed steps from copies of the same optimizer state, actual
online steps, and cumulative parameter changes from initialization.
Steps and cumulative changes use FP32 master weights when available,
with exported-checkpoint measurements reported separately. Their
sparsity need not equal gradient sparsity because momentum, clipping,
weight decay, and rounding can change which parameters move.

For the joint runs, we also compare which high-magnitude gradient
coordinates the teachers share, using Section~\ref{sec:subnetworks}'s
overlap measurement. Held-out per-domain capability, the macro-average,
and worst-domain change accompany the sparsity results. Curves show
generated-token exposure and GPU hours separately, with response counts,
updates, and teacher-scoring and diagnostic overhead recorded.

\subsection{Interpretation and decision rule}

At a fixed prefix, fresh unclipped sampled-token PG with a detached
log-ratio coefficient has the same expected gradient as full-vocabulary
reverse KL \citep{yu2026denseupdates,oh2026vopd}. This identity holds with
fixed prefix weights; weighting an original rollout token by its eventual
response length need not preserve it. Clipping and top-$k$ truncation
also change the comparison, and equality in expectation does not equate
sampled gradients, nonlinear optimizer steps, or training trajectories.

We report paired-seed effect estimates and confidence intervals against
a material difference margin specified before examining results.
One-seed runs remain exploratory; prompt resampling measures evaluation
uncertainty, not training-seed variation. Intervals entirely within the
margin support no material sparsity difference in the tested regime;
intervals spanning negligible and material effects are inconclusive.
Only a stable, meaningful gap motivates
the optional sampling-variance and truncation controls in
Appendix~\ref{app:optional_loss}.

\missing{PG/top-$64$ gradients and capability; full-vocabulary references;
optimizer steps and cumulative changes.}
